\section{Related Work}
\label{sec:related_work}

The field of parameter-space model consolidation has seen rapid growth due to the immense scale of modern deep neural networks. In this section, we situate our work within the landscape of model merging, curvature-aware parameter fusion, and test-time adaptation.

\subsection{Linear Model Merging and Task Arithmetic}
Initial efforts in model merging focused on simple linear weight interpolation between checkpoints. Linear Mode Connectivity (LMC) studies \cite{frankle2020linear, entezari2021role, wortsman2022model} demonstrated that models fine-tuned from the same pretrained initialization often lie in the same basin of the loss landscape, enabling low-loss interpolation. Moving beyond checkpoint averaging, \citet{ilharco2022editing} introduced \emph{Task Arithmetic}, which demonstrates that editing model capabilities is possible by adding or subtracting scaled task vectors (defined as the difference between the task-specific expert weights and the base model weights). While highly intuitive, Task Arithmetic employs uniform, static merging coefficients across all layers of the network. This ignores the highly varying semantic representation capacities of different layers and blocks, leading to catastrophic representation interference when multiple heterogeneous task vectors are added simultaneously.

\subsection{Curvature-Aware Parameter Fusion}
Recognizing that uniform scaling fails to account for varying parameter sensitivities, several researchers have incorporated loss-landscape geometry into parameter fusion. The pioneering work of \emph{Fisher Merging} \cite{matena2022merging} utilizes the diagonal of the Fisher Information Matrix (FIM) as a proxy for second-order loss sensitivity to weight the parameters of different task experts. Similarly, RegCalMerge \cite{regcalmerge} utilizes regularized calibration schemes alongside elastic spatial regularizers to mitigate task interference. 

However, all existing curvature-aware methods share a critical, limiting bottleneck: they rely strictly on a **diagonal approximation** of the Hessian or Fisher matrix. While essential to make computation and storage tractable over millions or billions of parameters, a diagonal approximation completely discards the cross-parameter second-order interactions (the off-diagonal Hessian entries). In modern deep networks, parameters are highly coupled, meaning that the sensitivity of a weight is strongly dependent on the states of surrounding weights. Discarding these cross-terms leads to sub-optimal merges. In contrast, our proposed ACM resolves this fundamental bottleneck. By projecting search directions onto the low-dimensional $K$-dimensional subspace of task vectors, we can compute and invert the \emph{full, non-diagonal, cross-parameter projected Hessian} with zero diagonal approximation.

\subsection{Test-Time Adaptation and the Overfitting-Optimizer Paradox}
To avoid manual scaling, state-of-the-art methods formulate model merging as an optimization problem. AdaMerging \cite{yang2023adamerging}, PolyMerge \cite{polymerge}, and layer-wise zero-order search methods \cite{sanyal2023sanity} optimize layer-specific coefficients at deployment time. Since ground-truth labels are generally unavailable at deployment, these methods rely on \emph{Test-Time Adaptation} (TTA) using unsupervised objectives, primarily Shannon entropy minimization of the merged model's predictions on a local test stream. 

While empirically functional in some scenarios, this paradigm exhibits major theoretical vulnerabilities. As highlighted by the \emph{Overfitting-Optimizer Paradox} \cite{regcalmerge}, optimizing parameters on tiny, unlabeled batches of test data is prone to transductive overfitting, where the model adjusts coefficients to minimize entropy on the local batch by collapsing representations, leading to poor generalization on unseen data. Additionally, test-time optimization is extremely slow, requiring hundreds of gradient-descent or evolutionary search steps, making it unsuitable for real-time edge deployment. Our ACM completely departs from this paradigm by solving for the optimal merging coefficients analytically in a single training-free calibration step, resolving both transductive overfitting and test-time computational overhead.
